%

\pol{\chapter{Stożkowe}}
\ita{\chapter{Coniche}}

\pol{Nazwy elipsa, parabola i hiperbola zostaną wprowadzone przez Apoloniusza i zaczerpnięte z wczesnej pitagorejskiej terminologii dotyczącej przykładania pól. Gdy Pitagorejczycy przykładają prostokąt do odcinka (czyli ustawiają jego podstawę wzdłuż odcinka tak, że jeden koniec podstawy pokrywał się z końcem odcinka), mówili o „ellipsie'', „paraboli'' lub „hyperboli'' w zależności od tego, czy podstawa prostokąta była krótsza od odcinka, dokładnie mu odpowiadała, czy też go przekraczała.}
\ita{I nomi ellisse, parabola e iperbole furono introdotti da Apollonio e derivano dall'antica terminologia pitagorica relativa all'applicazione delle aree. Quando i Pitagorici applicavano un rettangolo a un segmento (cioè disponevano la sua base lungo il segmento in modo che un estremo della base coincidesse con un estremo del segmento), parlavano di ``ellisse'', ``parabola'' o ``iperbole'' a seconda che la base del rettangolo fosse più corta del segmento, coincidesse esattamente con esso oppure lo superasse.}
% to jest Eves, s. 30, s. 31

\todofoot{225 BC - Apollonius of Perga writes On Conic Sections and names the ellipse, parabola, and hyperbola}

\pol{\section{Elipsa}}
\ita{\section{Ellisse}}

\pol{Tekst o elipsie...}
\ita{Testo sull'ellisse...}

\todofoot{Ogniska elipsy i hiperboli, ognisko, kierownica i mimośród stożkowych, asymptoty hiperboli, konstrukcja stycznej do stożkowej, rzuty ustalonego ogniska na styczne, własności izogonalne stożkowych, równania kanoniczne stożkowych, elipsa jako przekrój walca.}

\pol{\section{Parabola}}
\ita{\section{Parabola}}

\pol{Tekst o paraboli...}
\ita{Testo sulla parabola...}

\pol{Definicja paraboli - \cite[s. 247]{hartshorne2000}.}
\ita{Definizione di parabola - \cite[s. 247]{hartshorne2000}.}

\pol{\section{Hiperbola}}
\ita{\section{Iperbole}}

\pol{Tekst o hiperboli...}
\ita{Testo sull'iperbole...}

% TODO: https://bookstore.ams.org/browse?Author=%22A.%20V.%20Akopyan%22

% TODO:https://geometry.ru/books/conic_e.pdf

\todofoot{Ognisko, kierownica i mimośród stożkowej na przekroju stożka.}
\todofoot{Przekroje stożków ze sferami wpisanymi.}
\todofoot{Równanie ogólne stożkowej w układzie współrzędnych, duży i mały wyznacznik.}
\todofoot{Równania stożkowych we współrzędnych biegunowych.}

Audin \cite[s. 183]{audin_2003} % conics, quadrics
% https://en.wikipedia.org/wiki/Circumconic_and_inconic

% https://en.wikipedia.org/wiki/Five_points_determine_a_conic#Synthetic_proof
\todofoot{pięć punktów wyznacza stożkową}

% https://en.wikipedia.org/wiki/Parabola
\pol{Parabola: krzywa będąca zbiorem punktów równoodległych od prostej zwanej kierownicą paraboli i punktu zwanego ogniskiem paraboli.}
\ita{Parabola: curva costituita dall'insieme dei punti equidistanti da una retta detta direttrice della parabola e da un punto detto fuoco della parabola.}

% https://en.wikipedia.org/wiki/Hyperbola
% https://en.wikipedia.org/wiki/Conic_section + https://en.wikipedia.org/wiki/Degenerate_conic + https://en.wikipedia.org/wiki/Generalized_conic

% https://en.wikipedia.org/wiki/Dandelin_spheres
\pol{In geometry, the Dandelin spheres are one or two spheres that are tangent both to a plane and to a cone that intersects the plane. The intersection of the cone and the plane is a conic section, and the point at which either sphere touches the plane is a focus of the conic section, so the Dandelin spheres are also sometimes called focal spheres.[1]}
\ita{In geometria, le sfere di Dandelin sono una o due sfere tangenti sia a un piano sia a un cono che interseca tale piano. L'intersezione tra il cono e il piano è una sezione conica, e il punto in cui ciascuna sfera tocca il piano è un fuoco della conica; per questo motivo le sfere di Dandelin sono talvolta chiamate anche sfere focali.[1]}

\pol{The Dandelin spheres were discovered in 1822.[1][2] They are named in honor of the French mathematician Germinal Pierre Dandelin, though Adolphe Quetelet is sometimes given partial credit as well.[3][4][5]}
\ita{Le sfere di Dandelin furono scoperte nel 1822.[1][2] Prendono il nome dal matematico francese Germinal Pierre Dandelin, sebbene talvolta parte del merito venga attribuita anche ad Adolphe Quetelet.[3][4][5]}

\pol{The Dandelin spheres can be used to give elegant modern proofs of two classical theorems known to Apollonius. The first theorem is that a closed conic section (i.e. an ellipse) is the locus of points such that the sum of the distances to two fixed points (the foci) is constant. The second theorem is that for any conic section, the distance from a fixed point (the focus) is proportional to the distance from a fixed line (the directrix), the constant of proportionality being called the eccentricity.[6]}
\ita{Le sfere di Dandelin permettono di fornire eleganti dimostrazioni moderne di due teoremi classici già noti ad Apollonio. Il primo afferma che una conica chiusa (cioè un'ellisse) è il luogo dei punti per i quali la somma delle distanze da due punti fissi (i fuochi) è costante. Il secondo afferma che, per ogni sezione conica, la distanza da un punto fisso (il fuoco) è proporzionale alla distanza da una retta fissa (la direttrice); la costante di proporzionalità è detta eccentricità.[6]}

% https://www.deltami.edu.pl/2014/10/sprobuj-zostac-archimedesem/

% https://www.deltami.edu.pl/2024/12/zle-rysuje-stozki/

%